\sectioncentered*{Заключение}
\addcontentsline{toc}{section}{Заключение}

В лабораторной работе реализована универсальная модель релаксационной (сети Хопфилда) сети для задачи распознавания образов, покрывающая все четыре комбинации типа состояния (дискретное, непрерывное) и режима дискретного времени (синхронный, асинхронный). Веса рассчитываются методом проекции / дельта-проекции, функции активации~\mdash  модифицированная функция знака (дискретное состояние) и гиперболический тангенс (непрерывное состояние). Реализована проверка достижения состояния релаксации с детекцией зацикливания и подсчётом числа итераций.

Показано, что при чистой проекции ($\delta = 0$) и коэффициенте усиления $g = 2$ эталоны являются точными фиксированными точками сети: при $Q = 20$ эталонов, $n = 100$ успешность распознавания составляет $\ge 0.97$ для шума до $p = 0.2$ и остаётся высокой ($\ge 0.93$) при увеличении числа эталонов до $20$. Энергия монотонно убывает во всех режимах, синхронный режим сходится медленнее асинхронного.

\textbf{Семантическая составляющая.} Релаксационная сеть интерпретируется как динамическая система, спускающаяся по функции Ляпунова к локальным минимумам~\mdash  аттракторам, соответствующим запомненным образам; зашумлённый вход отображается в ближайший аттрактор, что и реализует распознавание.

\textbf{Прагматическая составляющая.} Получен воспроизводимый стенд на Python с интерактивным запуском, автоматической серией экспериментов (энергия, шум, ёмкость, визуальная сетка), сохранением метрик и весов, пригодный для сравнения всех режимов сети Хопфилда и подбора параметров $\delta$, $g$, $\varepsilon$.

\begin{thebibliography}{9}
	\addcontentsline{toc}{section}{Список использованных источников}

	\bibitem{method}
	Методические указания к лабораторной работе №4 по дисциплине <<Модели решения задач в интеллектуальных системах>>. \mdash  Минск : БГУИР.

	\bibitem{haykin}
	Хайкин С. Нейронные сети: полный курс. \mdash  2-е изд. \mdash  М. : Вильямс, 2006. \mdash  1104~с.

	\bibitem{numpy}
	NumPy documentation [Электронный ресурс]. \mdash  Режим доступа: \url{https://numpy.org/doc/}. \mdash  Дата доступа: 15.10.2026.

	\bibitem{matplotlib}
	Matplotlib documentation [Электронный ресурс]. \mdash  Режим доступа: \url{https://matplotlib.org/stable/}. \mdash  Дата доступа: 15.10.2026.
\end{thebibliography}
